\chapter{Localization of index and weak convergence of variations}\label{sec:index-convergence}
Negative directions persist only when their entire quadratic form remains negative definite. This section separates that finite-dimensional fact from regularity of a weak limit. Sources are \cite{Sharp17}, pp.~324--327, and the second-variation framework of the preceding chapters.

\begin{definition}[Compactly supported index]\label{def:compact-support-index}
For a smooth manifold $\Sigma$ and locally bounded real potential $q$, set $Q(f,h)=\int(\langle\nabla f,\nabla h\rangle-qfh)$ on compactly supported smooth interior functions. Its index is the supremum of the dimensions of vector spaces on which $Q$ is negative definite, allowing the value infinity. Stability means index zero. For a two-sided minimal hypersurface use $q=|B|^2+\operatorname{Ric}(\nu,\nu)$. This compact-support convention extends the index used in \cite{Sharp17}, Section 3, to noncompact domains.
\end{definition}

\begin{proposition}[Finite index gives stability outside a compact set]\label{prop:finite-index-stable-outside}
If the quadratic form in Definition \ref{def:compact-support-index} has finite index, there is a compact subset $K\subset\Sigma$ such that $Q(f,f)\geq0$ for every $f\in C_c^\infty(\Sigma\setminus K)$.
\end{proposition}
\begin{proof}
If no such $K$ existed, choose a negative test $f_1$, then a negative test $f_2$ outside a compact neighborhood of its support, and continue outside compact neighborhoods of all preceding supports. This produces arbitrarily many nonzero negative tests with pairwise disjoint supports. Every cross term vanishes, so their span is negative definite. They are linearly independent, contradicting finite index. This is the same disjoint-support mechanism as Lemma \ref{lem:disjoint-unstable-regions}; it uses no curvature estimate and asserts no geometry of the ends.
\end{proof}

\begin{proposition}[Finite index gives a stable punctured neighborhood]\label{prop:finite-index-stable-puncture}
Let $p$ be an interior point of a smooth manifold, and let $Q$ be a quadratic form as in Definition \ref{def:compact-support-index} on a punctured neighborhood $B_r(p)\setminus\{p\}$. If its compactly supported index there is finite, some smaller punctured ball is stable. The potential may be unbounded as $p$ is approached, provided it is locally bounded on the punctured domain.
\end{proposition}
\begin{proof}
Otherwise every smaller punctured ball admits a compactly supported negative test. After choosing such a test, its support stays a positive distance from $p$; choose the next ball with closure disjoint from that support. Induction produces arbitrarily many negative tests with disjoint supports in nested annuli. Their span is negative definite, contradicting finite index. This is the nested-annulus argument in \cite{Sharp17}, p.~326. It does not assert stability across $p$ when the potential or the underlying geometry is singular there.
\end{proof}

\begin{lemma}[Negative definite forms persist under matrix convergence]\label{lem:negative-form-persistence}
Let symmetric quadratic forms $Q_j$ on $\mathbb R^k$ converge coefficientwise to a negative definite form $Q$. Then for some $c>0$ and all sufficiently large $j$,
\[
 Q_j(z)\leq-c|z|^2\qquad(z\in\mathbb R^k).
\]
Having $Q(e_a)<0$ for each standard basis vector alone would not suffice.
\end{lemma}
\begin{proof}
The maximum of $Q$ on the compact unit sphere is strictly negative, say $-2c$. Coefficient convergence in finite dimension gives uniform convergence there, so eventually $Q_j\leq-c$ on the sphere. Homogeneity gives the assertion. For the final warning, the symmetric matrix with diagonal entries $-1$ and off-diagonal entries $2$ is negative on both coordinate vectors, but its value on $(1,1)$ is $2>0$. Thus one must control cross terms, not just each individual direction.
\end{proof}

\begin{lemma}[Weak continuity of a finite family of ambient second variations]\label{lem:weak-second-variation-continuity}
Let $M$ be smooth and closed, and let stationary $n$-varifolds $V_j\to V$ weakly. Fix smooth ambient vector fields $X_1,\ldots,X_k$, put $X_z=\sum_a z_aX_a$, and define
\[
 \varphi_t^z(x)=\exp_x(tX_z(x)),\qquad
 Q_{V_j}(z)=\left.\frac{d^2}{dt^2}\right|_{t=0}\|(\varphi_t^z)_\#V_j\|(M).
\]
These symmetric quadratic forms converge coefficientwise to $Q_V$. For a smooth minimal hypersurface, the form is the usual second variation of its normal velocity. Compare \cite{Sharp17}, p.~324.
\end{lemma}
\begin{proof}
For $z$ in a fixed bounded set and $t$ sufficiently small, the maps are smooth diffeomorphisms close to the identity. The area formula gives their masses by integrating $J_P\varphi_t^z(x)$ over $G_n(M)$. Its second time derivative at zero is a continuous function of $(x,P)$ depending quadratically on $z$: it is the Hessian at zero of the smooth finite-parameter family $x\mapsto\exp_x(\sum_a s_aX_a(x))$, evaluated on $(z,z)$. All its derivatives in the finite parameters are uniformly bounded on the compact Grassmann bundle near zero, so differentiation under the integral is valid.
Weak convergence makes the integral of each coefficient converge. The constant test function also gives uniform mass bounds. Stationarity of $V$ follows from weak continuity of first variation. For a smooth minimal hypersurface, tangential velocity is infinitesimal reparametrization and the acceleration term pairs with zero mean curvature; the second-variation formula therefore identifies the resulting form with the Jacobi quadratic form of $X_z^\perp$. No bound on the second fundamental forms of the approximating hypersurfaces was used in the weak-continuity argument.
\end{proof}

\begin{theorem}[Index cannot increase at a smooth weak limit]\label{thm:index-smooth-varifold-limit}
Let $M^{n+1}$ be smooth and closed, $n\geq1$. Let $\Sigma_j$ and $\Sigma$ be smooth closed two-sided embedded minimal hypersurfaces. If
\[
 |\Sigma_j|\longrightarrow m|\Sigma|\quad\text{as varifolds},\qquad m>0,
\]
then $\operatorname{index}(\Sigma)\leq\liminf_j\operatorname{index}(\Sigma_j)$. The smoothness and two-sidedness of the limit are assumptions, not conclusions. This proves the index-persistence step underlying \cite{Sharp17}, Section 4, Claim 3, in this setting.
\end{theorem}
\begin{proof}
Take any finite-dimensional negative subspace for the Jacobi form of $\Sigma$, with smooth basis $f_1,\ldots,f_k$. Extend $f_a\nu$ to smooth ambient fields $X_a$ using a tubular neighborhood and a cutoff. At the limit varifold, Lemma \ref{lem:weak-second-variation-continuity} identifies the form on coefficient vectors with
\[
 Q_{m|\Sigma|}(z)=mQ_\Sigma\left(\sum_a z_af_a,\sum_a z_af_a\right),
\]
which is negative definite because $m>0$. Lemma \ref{lem:negative-form-persistence} makes the corresponding forms on $\Sigma_j$ negative definite for all large $j$. Their normal velocities are the functions $\langle X_a,\nu_j\rangle$. These functions are linearly independent: a zero linear combination would have second variation zero, contradicting strict negativity for its nonzero coefficient vector. Their span therefore gives $\operatorname{index}(\Sigma_j)\geq k$. Since every finite negative subspace of the limit can be used, the stated inequality follows, including when its index is infinite. This argument neither establishes smooth graphical convergence nor removes singularities of a weak limit.
\end{proof}

\begin{corollary}[Stability passes to a smooth positive-multiplicity limit]\label{cor:stable-smooth-varifold-limit}
Under the hypotheses of Theorem \ref{thm:index-smooth-varifold-limit}, if every $\Sigma_j$ is stable, then $\Sigma$ is stable, irrespective of the positive multiplicity $m$.
\end{corollary}
\begin{proof}
The theorem gives $\operatorname{index}(\Sigma)\leq0$, which is stability. Equivalently, the second variation of any ambient extension of a normal test on $\Sigma$ is the limit of nonnegative second variations, and the positive factor $m$ can be divided out. If the limiting support is singular, this corollary does not assert that it is a smooth stable hypersurface.
\end{proof}
